% Options for packages loaded elsewhere
\PassOptionsToPackage{unicode}{hyperref}
\PassOptionsToPackage{hyphens}{url}
\documentclass[
]{ltjsarticle}
\usepackage{xcolor}
\usepackage{amsmath,amssymb}
\setcounter{secnumdepth}{-\maxdimen} % remove section numbering
\usepackage{iftex}
\ifPDFTeX
  \usepackage[T1]{fontenc}
  \usepackage[utf8]{inputenc}
  \usepackage{textcomp} % provide euro and other symbols
\else % if luatex or xetex
  \usepackage{unicode-math} % this also loads fontspec
  \defaultfontfeatures{Scale=MatchLowercase}
  \defaultfontfeatures[\rmfamily]{Ligatures=TeX,Scale=1}
\fi
\usepackage{lmodern}
\ifPDFTeX\else
  % xetex/luatex font selection
\fi
% Use upquote if available, for straight quotes in verbatim environments
\IfFileExists{upquote.sty}{\usepackage{upquote}}{}
\IfFileExists{microtype.sty}{% use microtype if available
  \usepackage[]{microtype}
  \UseMicrotypeSet[protrusion]{basicmath} % disable protrusion for tt fonts
}{}
\makeatletter
\@ifundefined{KOMAClassName}{% if non-KOMA class
  \IfFileExists{parskip.sty}{%
    \usepackage{parskip}
  }{% else
    \setlength{\parindent}{0pt}
    \setlength{\parskip}{6pt plus 2pt minus 1pt}}
}{% if KOMA class
  \KOMAoptions{parskip=half}}
\makeatother
\usepackage{longtable,booktabs,array}
\newcounter{none} % for unnumbered tables
\usepackage{calc} % for calculating minipage widths
% Correct order of tables after \paragraph or \subparagraph
\usepackage{etoolbox}
\makeatletter
\patchcmd\longtable{\par}{\if@noskipsec\mbox{}\fi\par}{}{}
\makeatother
% Allow footnotes in longtable head/foot
\IfFileExists{footnotehyper.sty}{\usepackage{footnotehyper}}{\usepackage{footnote}}
\makesavenoteenv{longtable}
\setlength{\emergencystretch}{3em} % prevent overfull lines
\providecommand{\tightlist}{%
  \setlength{\itemsep}{0pt}\setlength{\parskip}{0pt}}
\usepackage{bookmark}
\IfFileExists{xurl.sty}{\usepackage{xurl}}{} % add URL line breaks if available
\urlstyle{same}
\hypersetup{
  hidelinks,
  pdfcreator={LaTeX via pandoc}}

\author{}
\date{}

\begin{document}

\section{信号論・制御理論と量子力学・量子光学・開放量子系の構造的対応関係に関する観察}\label{ux4fe1ux53f7ux8ad6ux5236ux5fa1ux7406ux8ad6ux3068ux91cfux5b50ux529bux5b66ux91cfux5b50ux5149ux5b66ux958bux653eux91cfux5b50ux7cfbux306eux69cbux9020ux7684ux5bfeux5fdcux95a2ux4fc2ux306bux95a2ux3059ux308bux89b3ux5bdf}

\textbf{著者}：木原 範昭（Noriaki Kihara） \textbf{所属}：WF System Co.,
Ltd.~/ 大阪大学基礎工学部（卒業）
\textbf{ORCID}：\href{https://orcid.org/0009-0004-6753-4020}{0009-0004-6753-4020}
\textbf{版}：v6 \textbf{日付}：2026 年 6 月 3 日 \textbf{ライセンス}：CC
BY 4.0 \textbf{Concept DOI}：（v1 公開時に付与） \textbf{v1 Version
DOI}：（v1 公開時に付与） \textbf{v2 Version DOI}：（v2 公開時に付与）
\textbf{v3 Version DOI}：（v3 公開時に付与） \textbf{v4 Version
DOI}：（v4 公開時に付与） \textbf{v5 Version DOI}：（v5 公開時に付与）
\textbf{v6 Version DOI}：（v6 公開時に付与）

\begin{center}\rule{0.5\linewidth}{0.5pt}\end{center}

\subsection{本稿の性格}\label{ux672cux7a3fux306eux6027ux683c}

\textbf{本稿は観察論文である。証明論文・主張論文ではない。}

本稿は新たな物理的予言を行わない。新たな数学的定理を証明しない。新たな解釈を提案しない。

本稿が行うのは、\textbf{信号論・制御理論において既に確立され社会実装されている数学的構造と、量子力学・量子光学・開放量子系において確立されている数学的構造の間に、すでに
10 個の構造的対応関係が文献上存在することを並置して観察する}に留まる。

これらの対応関係は、両領域の標準教科書および古典原論文に既に記載されている事実である。本稿はそれらを横断的に整理するに過ぎない。

本稿は以下を主張しない：

\begin{itemize}
\tightlist
\item
  量子力学の新解釈を提案する
\item
  信号論・制御理論の物理学的拡張を主張する
\item
  いずれかの分野の優位性を主張する
\item
  物理定数の予言を行う
\item
  既存の物理理論の修正を提案する
\item
  新たな数学的定理を主張する
\end{itemize}

評価・解釈は読者に委ねる。

\subsubsection{v2
での重要な用語選択}\label{v2-ux3067ux306eux91cdux8981ux306aux7528ux8a9eux9078ux629e}

v1 では「構造的同一性 (structural identity)」を用いたが、v2
では「\textbf{構造的対応関係 (structural
correspondence)}」と表現する。これは、§2 で列挙する 10
項目のうち、いくつかは厳密な数学的同一性（identity）であるが、いくつかは数学的構造の部分的並行性（analogy）に留まるためである。両者の峻別は
§2 冒頭の分類表に示す。

\begin{center}\rule{0.5\linewidth}{0.5pt}\end{center}

\subsection{要旨}\label{ux8981ux65e8}

理論物理学（量子力学・量子光学・開放量子系）と工学（信号論・制御理論）は
20
世紀において独立に発展した。前者は微視的物理現象の記述を目的とし、後者は通信・計測・制御の実用化を目的とした。

しかし両領域の根底にある数学的構造は、いくつかの重要な点で厳密な同一性または強い構造的対応を持つ。特に、Heisenberg
の量子力学的不確定性
\(\Delta x \cdot \Delta p \geq \hbar/2\)（Heisenberg 1927）と、Gabor
の時間-帯域不確定性 \(\Delta t \cdot \Delta \omega \geq 1/2\)（Gabor
1946）は同一の数学的定理であり、Gabor 自身が両者の等価性を明示している。

本稿は、この既知の対応関係を起点として、両領域の構造的対応関係を 10
項目に整理して観察する。各項目について、両領域の標準文献を併記し、厳密な同一性が成立する範囲と、構造的並行性に留まる範囲を明示する。

特筆すべき事実として、(i)
同一の数学的対象が両領域で独立に発見されている（Wigner 1932
が量子力学、Ville 1948 が信号論で同等の位相空間擬確率分布）、(ii)
Kennard-Robertson 型の Fourier 不確定性と Gabor
の時間-帯域不確定性が同じ Fourier 解析的不等式に基づくこと、および Gabor
自身による対応関係の明示、(iii)
信号論・制御理論側の数学構造はすでに通信・GPS・自動運転・レーダー・MRI・光通信等の工学領域で社会実装され日常的に動作している事実、を観察する。

\begin{center}\rule{0.5\linewidth}{0.5pt}\end{center}

\subsection{§1 序論}\label{ux5e8fux8ad6}

\subsubsection{1.1
二領域の独立発展}\label{ux4e8cux9818ux57dfux306eux72ecux7acbux767aux5c55}

20 世紀前半、量子力学は Heisenberg、Schrödinger、Dirac、von Neumann
らによって体系化された {[}1, 11, 16, 17{]}。同時期、信号論・制御理論は
Nyquist、Shannon、Wiener、Gabor、Kalman らによって体系化された {[}2, 5,
6, 7, 8{]}。

両領域は
\textbf{目的・動機・適用範囲を異にしながら独立に発展}した。前者は微視的物理現象の予言と理解を目的とし、後者は通信・計測・制御における実用問題の解決を目的とした。

\subsubsection{1.2
数学的構造の対応関係}\label{ux6570ux5b66ux7684ux69cbux9020ux306eux5bfeux5fdcux95a2ux4fc2}

しかし、両領域の数学的構造を並置すると、複数の重要な概念について
\textbf{厳密な同一性または強い構造的対応}
が成立していることが確認できる。本稿はこの観察を整理する。

最も明示的な例は、Heisenberg の量子力学的不確定性原理 {[}1{]} と、Gabor
の信号論的不確定性原理 {[}2{]} の同一性である。Gabor (1946) {[}2{]}
は自身の論文において、両者が同一の数学的構造を持つことを明示的に指摘している。

\subsubsection{1.3 本稿の方法}\label{ux672cux7a3fux306eux65b9ux6cd5}

本稿は以下の方法を採る：

\begin{enumerate}
\def\labelenumi{\arabic{enumi}.}
\tightlist
\item
  両領域から、対応関係を持つと観察される概念ペアを選定する
\item
  各ペアについて、両領域の標準文献（古典原論文または定本教科書）を併記する
\item
  数学的構造の対応関係を表形式で並置する
\item
  \textbf{厳密な同一性が成立する範囲と、構造的並行性に留まる範囲を明示する}
\item
  解釈・評価は加えない
\end{enumerate}

本稿は両領域の専門家にとって自明な事実の整理に留まり、新規性は構造の「並置と分類」のみにある。新規な数学的・物理的主張は一切行わない。

\subsubsection{1.4 本稿の構成}\label{ux672cux7a3fux306eux69cbux6210}

§2 で 10
個の構造的対応関係を順次列挙する。各項目で「厳密に同一な部分」と「構造的並行性に留まる部分」を区別する。§3
で観察事項のまとめを行う。§4 で本稿が主張しないことを明示する。§5
で含意を列挙する（これらも主張ではなく観察である）。§6 で結語を述べる。

\begin{center}\rule{0.5\linewidth}{0.5pt}\end{center}

\subsection{§2 10
個の構造的対応関係}\label{ux500bux306eux69cbux9020ux7684ux5bfeux5fdcux95a2ux4fc2}

\subsubsection{2.0
対応関係の分類}\label{ux5bfeux5fdcux95a2ux4fc2ux306eux5206ux985e}

以下の 10 項目を、対応関係の強さによって分類しておく。

{\def\LTcaptype{none} % do not increment counter
\begin{longtable}[]{@{}
  >{\raggedright\arraybackslash}p{(\linewidth - 4\tabcolsep) * \real{0.3333}}
  >{\raggedright\arraybackslash}p{(\linewidth - 4\tabcolsep) * \real{0.3333}}
  >{\raggedright\arraybackslash}p{(\linewidth - 4\tabcolsep) * \real{0.3333}}@{}}
\toprule\noalign{}
\begin{minipage}[b]{\linewidth}\raggedright
\#
\end{minipage} & \begin{minipage}[b]{\linewidth}\raggedright
概念ペア
\end{minipage} & \begin{minipage}[b]{\linewidth}\raggedright
分類
\end{minipage} \\
\midrule\noalign{}
\endhead
\bottomrule\noalign{}
\endlastfoot
1 & 不確定性原理 &
\textbf{厳密な数学的同一性}（単位スケール除き同一定理） \\
2 & 時間-周波数擬確率分布 & \textbf{厳密な数学的同一性}（定義式同一） \\
3 & パラキシャル方程式 / Schrödinger 方程式 &
\textbf{条件付き厳密同型}（近軸・単色・スカラー近似下の PDE
形式として同一） \\
4 & サンプリング / 位相空間自由度 &
\textbf{強い構造的対応}（有限領域における有効自由度カウントの形式的対応、実効自由度の漸近的同一） \\
5 & 状態空間表現 / Hilbert 空間描像 &
\textbf{構造的並行性}（線形時間発展だが、ユニタリ性・入出力構造に相違） \\
6 & 可観測性 / CSCO &
\textbf{構造的並行性}（並行する形式問題、ただし動力学／運動学の機構差） \\
7 & Kalman フィルタ / 量子フィルタリング & \textbf{構造的並行性}（POVM
経由で対応、射影測定とは別物） \\
8 & Jones ベクトル / qubit & \textbf{厳密な数学的同一性}（同じ 2
成分複素ベクトル、\(U(2)\) 作用） \\
9 & デフェージング / 位相雑音 & \textbf{構造的並行性}（T₂
過程の数学のみ共有、デコヒーレンス全体ではない） \\
10 & SVD / Schmidt 分解 &
\textbf{厳密な数学的同一性}（線形代数のツール同一） \\
\end{longtable}
}

厳密な数学的同一性が成立するのは \#1, \#2, \#3, \#8, \#10 の 5
項目、強い構造的対応は \#4 の 1 項目、構造的並行性に留まるのは \#5, \#6,
\#7, \#9 の 4 項目である。

\subsubsection{2.1
不確定性原理（厳密な数学的同一性）}\label{ux4e0dux78baux5b9aux6027ux539fux7406ux53b3ux5bc6ux306aux6570ux5b66ux7684ux540cux4e00ux6027}

{\def\LTcaptype{none} % do not increment counter
\begin{longtable}[]{@{}
  >{\raggedright\arraybackslash}p{(\linewidth - 2\tabcolsep) * \real{0.5000}}
  >{\raggedright\arraybackslash}p{(\linewidth - 2\tabcolsep) * \real{0.5000}}@{}}
\toprule\noalign{}
\begin{minipage}[b]{\linewidth}\raggedright
信号論
\end{minipage} & \begin{minipage}[b]{\linewidth}\raggedright
量子力学
\end{minipage} \\
\midrule\noalign{}
\endhead
\bottomrule\noalign{}
\endlastfoot
\textbf{Gabor 時間-帯域不確定性} & \textbf{Heisenberg-Kennard-Robertson
不確定性} \\
\(\Delta t \cdot \Delta \omega \geq \dfrac{1}{2}\) &
\(\Delta x \cdot \Delta p \geq \dfrac{\hbar}{2}\) \\
{[}Gabor 1946{]} {[}2{]} & {[}Heisenberg 1927{]} {[}1{]}, {[}Kennard
1927{]} {[}22{]}, {[}Robertson 1929{]} {[}23{]} \\
\end{longtable}
}

両者は単位（\(\hbar = 1\)
の自然単位）を除いて同一の不等式である。両者の数学的証明はいずれも
Cauchy-Schwarz 不等式と Fourier
変換のユニタリ性に基づき、構造的に同一である。位置 \(x\) 表示の Fourier
共役が運動量 \(p/\hbar = k\)
である事実から、\(\sigma_x \sigma_p \geq \hbar/2\) は Gabor
不等式と同一の Fourier 解析的命題となる。

\textbf{歴史的注記}：Heisenberg (1927) {[}1{]}
は不確定性原理を直観的・操作的議論として提示した。これを現代的な分散不等式
\(\sigma_x \sigma_p \geq \hbar/2\) の厳密形に定式化したのは Kennard
(1927) {[}22{]} であり、任意の自己共役演算子対への一般化を行ったのは
Robertson (1929) {[}23{]} である。本稿で言う「Heisenberg
不確定性」はこれら 3 論文に基づく現代的標準形を指す。

Gabor (1946) {[}2{]} は自著において、自らの不等式が「波動力学における
Heisenberg
の不確定性原理の類似（analog）である」と明示的に述べている。すなわち\textbf{両者の等価性は信号論側の創始者自身によって
1946 年の時点で言明されている}。

\subsubsection{2.2
時間-周波数（位相空間）擬確率分布（厳密な数学的同一性）}\label{ux6642ux9593-ux5468ux6ce2ux6570ux4f4dux76f8ux7a7aux9593ux64ecux78baux7387ux5206ux5e03ux53b3ux5bc6ux306aux6570ux5b66ux7684ux540cux4e00ux6027}

{\def\LTcaptype{none} % do not increment counter
\begin{longtable}[]{@{}
  >{\raggedright\arraybackslash}p{(\linewidth - 2\tabcolsep) * \real{0.5000}}
  >{\raggedright\arraybackslash}p{(\linewidth - 2\tabcolsep) * \real{0.5000}}@{}}
\toprule\noalign{}
\begin{minipage}[b]{\linewidth}\raggedright
信号論
\end{minipage} & \begin{minipage}[b]{\linewidth}\raggedright
量子力学
\end{minipage} \\
\midrule\noalign{}
\endhead
\bottomrule\noalign{}
\endlastfoot
\textbf{Wigner-Ville 分布} & \textbf{Wigner 関数} \\
\(W(t,\omega) = \int s(t+\tau/2)\,\overline{s(t-\tau/2)}\, e^{-i\omega\tau}\, d\tau\)
&
\(W(x,p) = \dfrac{1}{2\pi\hbar}\int \psi^*(x+y/2)\,\psi(x-y/2)\, e^{ipy/\hbar}\, dy\) \\
{[}Ville 1948{]} {[}4{]} & {[}Wigner 1932{]} {[}3{]} \\
\end{longtable}
}

両者は同一の数学的定義に従う擬確率分布である。\textbf{Wigner（量子力学、1932
年 {[}3{]}）と Ville（信号論、1948 年
{[}4{]}）は別人}であるが、両者が独立に発見した数学的対象は同一の構造を持つ。両分野の標準名称が「\textbf{Wigner-Ville
分布}」と二名を冠することに、独立発見と数学的同一性の歴史が反映されている。

両分布とも以下の同一性質を持つ：

\begin{itemize}
\tightlist
\item
  周辺分布が真の確率分布（位置分布・運動量分布／時間分布・周波数分布）を与える
\item
  負値を取りうる（擬確率分布、正値性は満たさない）
\item
  線形 Fourier 変換に対する共変性
\item
  同じ自己相関構造を持つ
\end{itemize}

両者は信号処理（レーダー、ソナー、時間-周波数解析）および量子光学（quantum
tomography）において \textbf{同一の数学的対象として}
日常的に使用されている。

\subsubsection{2.3 Schrödinger
方程式とパラキシャル波動方程式（条件付き厳密同型：近軸近似下の PDE
形式）}\label{schruxf6dinger-ux65b9ux7a0bux5f0fux3068ux30d1ux30e9ux30adux30b7ux30e3ux30ebux6ce2ux52d5ux65b9ux7a0bux5f0fux6761ux4ef6ux4ed8ux304dux53b3ux5bc6ux540cux578bux8fd1ux8ef8ux8fd1ux4f3cux4e0bux306e-pde-ux5f62ux5f0f}

{\def\LTcaptype{none} % do not increment counter
\begin{longtable}[]{@{}
  >{\raggedright\arraybackslash}p{(\linewidth - 2\tabcolsep) * \real{0.5000}}
  >{\raggedright\arraybackslash}p{(\linewidth - 2\tabcolsep) * \real{0.5000}}@{}}
\toprule\noalign{}
\begin{minipage}[b]{\linewidth}\raggedright
信号論（光学）
\end{minipage} & \begin{minipage}[b]{\linewidth}\raggedright
量子力学
\end{minipage} \\
\midrule\noalign{}
\endhead
\bottomrule\noalign{}
\endlastfoot
\textbf{パラキシャル波動方程式} & \textbf{Schrödinger 方程式} \\
\(i\dfrac{\partial E}{\partial z} = -\dfrac{1}{2k}\nabla_\perp^2 E - \dfrac{k\,\delta n}{n_0}E\)
&
\(i\hbar \dfrac{\partial \psi}{\partial t} = -\dfrac{\hbar^2}{2m}\nabla^2 \psi + V\psi\) \\
光ファイバー、レーザー伝播、ホログラフィー &
単一粒子の非相対論的量子力学 \\
{[}Saleh \& Teich 1991{]} {[}12{]} & {[}Schrödinger 1926{]} {[}16{]} \\
\end{longtable}
}

両者は \textbf{近軸・単色・スカラー近似のもとで}、変数の対応
\((z \leftrightarrow t,\ 1/k \leftrightarrow \hbar/m,\ V/\hbar \leftrightarrow -\,k\delta n/n_0)\)
の下で \textbf{数学的に同型の偏微分方程式}
となる。ここで標準的な搬送波規約 \(E \propto e^{ikz}\)
を採ると屈折率項は負符号で現れ、屈折率の高い領域（\(\delta n>0\)、導波）が\textbf{束縛ポテンシャル井戸（\(V<0\)）}に対応する（相対マイナス符号）。なお斥力ポテンシャル（\(V>0\)）も正当な
Schrödinger
型であり、この符号は同型分類を損なわず、明示辞書の整合のための注記である。

この同型性は光ファイバー通信および集積光学の設計において日常的に利用されており、Saleh
\& Teich (1991) {[}12{]} 第 2-3
章を含む光学標準教科書に明示されている。光学側では「\textbf{量子力学の数学を光学に流用する}」と教えられることもあれば、逆に「\textbf{光学の数学を量子力学に流用する}」と教えられることもあるが、いずれの方向でも数学的内容は同じである。

\textbf{注意}：両者の対応はあくまで近似的・数学的同型である。光学側の
\(z\) は伝播距離であり、量子力学側の \(t\)
は時間である。物理的意味（確率解釈、Hilbert
空間内積構造、測定理論など）まで同一となるわけではなく、対応関係は線形
Schrödinger 型 PDE の数学的形式にとどまる。

\subsubsection{2.4
サンプリング定理と位相空間自由度（強い構造的対応）}\label{ux30b5ux30f3ux30d7ux30eaux30f3ux30b0ux5b9aux7406ux3068ux4f4dux76f8ux7a7aux9593ux81eaux7531ux5ea6ux5f37ux3044ux69cbux9020ux7684ux5bfeux5fdc}

{\def\LTcaptype{none} % do not increment counter
\begin{longtable}[]{@{}
  >{\raggedright\arraybackslash}p{(\linewidth - 2\tabcolsep) * \real{0.5000}}
  >{\raggedright\arraybackslash}p{(\linewidth - 2\tabcolsep) * \real{0.5000}}@{}}
\toprule\noalign{}
\begin{minipage}[b]{\linewidth}\raggedright
信号論
\end{minipage} & \begin{minipage}[b]{\linewidth}\raggedright
量子力学
\end{minipage} \\
\midrule\noalign{}
\endhead
\bottomrule\noalign{}
\endlastfoot
\textbf{Shannon-Nyquist サンプリング定理 / time-bandwidth product} &
\textbf{位相空間自由度カウント} \\
帯域 \(B\)、観測時間 \(T\) の実信号は、有効自由度 \(N \sim 2BT\)
で近似的に記述される（複素 baseband なら \(\sim BT\)
個の独立\textbf{複素}標本相当） & 位相空間体積 \(V \cdot p_{\max}^d\)
に対する量子状態数 \(\sim V \cdot p_{\max}^d / (2\pi\hbar)^d\) \\
{[}Nyquist 1928{]} {[}5{]}, {[}Shannon 1949{]} {[}6{]},
{[}Slepian-Pollak 1961-1964{]} {[}13{]} & {[}Planck 1906{]} {[}24{]},
{[}Sackur 1911; Tetrode 1912{]} {[}25{]}, 統計力学標準教科書 \\
\end{longtable}
}

両者は「\textbf{位相空間最小セル}」概念という共通の枠組みに基づく
\textbf{強い構造的対応} を持つ。信号論では帯域時間積（time-bandwidth
product）として、量子力学では位相空間体積を Planck 単位 \(h\)
で割った値として現れる。

\textbf{自由度カウントの形式は強く対応している}
が、両者は厳密に同一の定理ではない。Shannon-Nyquist／Slepian-Pollak
型の時間--帯域自由度と、半古典的状態数（Weyl law
による量子状態数）は、それぞれ異なる測度・境界条件・物理的解釈を持つ。Gabor
(1946) {[}2{]}
は信号論側でこの自由度カウントを「\textbf{logon}」と呼んで導入しており、量子力学側の位相空間セル概念と本質的に並行する数え上げである。

\textbf{精緻化}：信号論側では、有限時間・有限帯域の信号空間における有効自由度の厳密な定式化は
Slepian, Pollak, Landau による prolate spheroidal wave functions の研究
{[}13{]}
で確立されている。\textbf{時間制限かつ帯域制限の非零信号は厳密には存在しないため、信号空間は形式上無限次元であり、\(N \sim 2BT\)
は「集中固有値（prolate eigenvalues）が 1
に近い実効次元」の漸近的表現である。} 量子力学側では、位相空間体積を
\(h^d\) で割る数え上げは黒体放射の Planck 量子化（1906）{[}24{]}
や統計力学の Sackur-Tetrode 方程式（Sackur 1911; Tetrode 1912）{[}25{]}
に源流があり、半古典的状態数の Weyl law として一般化される。

情報理論的に表現すれば「\textbf{有限帯域・有限観測時間の系の独立自由度は実効的に有限}」であり、これが量子力学の
Hilbert
空間次元の有限化（位相空間体積による）と同じ形式の枠組みである。両者の対応は
\textbf{同一定理ではなく、同形式の自由度カウント}
という強い構造的対応である。

\subsubsection{2.5 状態空間表現と Hilbert
空間描像（構造的並行性、ただし重要な相違あり）}\label{ux72b6ux614bux7a7aux9593ux8868ux73feux3068-hilbert-ux7a7aux9593ux63cfux50cfux69cbux9020ux7684ux4e26ux884cux6027ux305fux3060ux3057ux91cdux8981ux306aux76f8ux9055ux3042ux308a}

{\def\LTcaptype{none} % do not increment counter
\begin{longtable}[]{@{}
  >{\raggedright\arraybackslash}p{(\linewidth - 2\tabcolsep) * \real{0.5000}}
  >{\raggedright\arraybackslash}p{(\linewidth - 2\tabcolsep) * \real{0.5000}}@{}}
\toprule\noalign{}
\begin{minipage}[b]{\linewidth}\raggedright
制御理論
\end{minipage} & \begin{minipage}[b]{\linewidth}\raggedright
量子力学
\end{minipage} \\
\midrule\noalign{}
\endhead
\bottomrule\noalign{}
\endlastfoot
\textbf{線形状態方程式} & \textbf{Schrödinger 方程式（線形版）} \\
状態方程式:
\(\dot{\mathbf{x}}(t) = A\,\mathbf{x}(t) + B\,\mathbf{u}(t)\) &
状態発展:
\(i\hbar\,\dfrac{d|\psi\rangle}{dt} = \hat{H}\,|\psi\rangle\) \\
出力方程式: \(\mathbf{y}(t) = C\,\mathbf{x}(t)\) & 観測値:
\(\langle A\rangle = \langle\psi|\hat{A}|\psi\rangle\) \\
{[}Kalman 1960{]} {[}8{]} & {[}von Neumann 1932{]} {[}11{]} \\
\end{longtable}
}

両者ともに \textbf{線形作用素（\(A\) または
\(\hat{H}\)）によって線形空間上での時間発展を記述する}
点で並行している。ただし、以下の重要な相違が存在する：

\textbf{相違 (i) ユニタリ性の有無}：量子力学では \(\hat{H}\)
がエルミート（自己共役）であるため \(e^{-i\hat{H}t/\hbar}\)
はユニタリ作用素となり、状態のノルム（確率）が保存される。一方、制御理論の
\(A\)
は一般に反エルミート（純虚スペクトル）である必要がなく、\textbf{実部を持つ固有値（不安定極・散逸モード）を許容する}。すなわち閉じた量子系の保存的時間発展と、開放系・散逸系を扱う制御系では、対応するのは
\(A\) が反エルミートかつ \(B\mathbf{u} = 0\) の特殊な場合のみである。

\textbf{相違 (ii) 観測値の双線形性}：量子力学の期待値
\(\langle A\rangle = \langle\psi|\hat{A}|\psi\rangle\) は状態
\(|\psi\rangle\) について \textbf{双線形（二次形式）}
である。一方、制御理論の出力 \(\mathbf{y} = C\mathbf{x}\) は状態
\(\mathbf{x}\) について \textbf{線形}
である。観測量の抽出方式が両者で構造的に異なる。

\textbf{相違 (iii) 駆動入力の有無}：制御理論の駆動入力
\(B\mathbf{u}(t)\) は被制御系の本質をなすが、閉じた Schrödinger
方程式にはこれに対応する自由項がない。開放量子系（Lindblad 方程式
{[}15{]}）まで拡張すると駆動項が加わるが、本対応は標準的な閉じた
Schrödinger 方程式に限った話である。

したがって、両者の対応は
\textbf{「線形空間上の線形時間発展」というレベルでの構造的並行性}
であり、有限次元量子系（qubit, qudit）に限ったとしても
\textbf{ユニタリ性・観測量抽出の方式・駆動入力の取り扱いに本質的相違が残る}。これらの相違は両領域の文献において明示されており、構造的対応の限界を示すものである。

\subsubsection{2.6
可観測性と状態識別可能性（構造的並行性、機構の相違あり）}\label{ux53efux89b3ux6e2cux6027ux3068ux72b6ux614bux8b58ux5225ux53efux80fdux6027ux69cbux9020ux7684ux4e26ux884cux6027ux6a5fux69cbux306eux76f8ux9055ux3042ux308a}

{\def\LTcaptype{none} % do not increment counter
\begin{longtable}[]{@{}
  >{\raggedright\arraybackslash}p{(\linewidth - 2\tabcolsep) * \real{0.5000}}
  >{\raggedright\arraybackslash}p{(\linewidth - 2\tabcolsep) * \real{0.5000}}@{}}
\toprule\noalign{}
\begin{minipage}[b]{\linewidth}\raggedright
制御理論
\end{minipage} & \begin{minipage}[b]{\linewidth}\raggedright
量子力学
\end{minipage} \\
\midrule\noalign{}
\endhead
\bottomrule\noalign{}
\endlastfoot
\textbf{観測可能性 (Observability)} & \textbf{完全観測可能集合
(CSCO)} \\
観測可能性行列
\(\mathcal{O} = \begin{pmatrix}C\\CA\\CA^2\\\vdots\\CA^{n-1}\end{pmatrix}\)
のランク条件 & 互いに可換な自己共役演算子の最大集合（Complete Set of
Commuting Observables） \\
状態が出力時系列から復元可能か &
状態が観測量の同時固有値から特定可能か \\
{[}Kalman 1960{]} {[}8{]} & {[}Dirac 1930{]} {[}17{]}, {[}von Neumann
1932{]} {[}11{]} \\
\end{longtable}
}

両者は「\textbf{観測から状態を特定するための条件}」に関する
\textbf{並行する形式化} である。ただし、機構には重要な相違がある：

\textbf{相違 (i) 動力学 vs 運動学}：Kalman
の観測可能性条件は、出力時系列 \(\{y(t)\}_{t \geq 0}\) から初期状態
\(\mathbf{x}(0)\) を一意に決定できる条件であり、行列 \(CA^k\)
の積み上げという \textbf{動力学} を利用する。一方、量子力学の CSCO
は、単一時刻における可換観測量代数が状態空間の基底を一意に指定する
\textbf{運動学的} 条件であり、時間発展を必要としない。

\textbf{相違 (ii) 単一軌道 vs 多数コピー}：Kalman
可観測性は決定論的な単一軌道から状態を再構成する条件である。一方、量子力学では単一の測定で一般に状態は決定できず、状態決定には
\textbf{量子状態トモグラフィ}（同一状態の多数コピーに対する繰り返し測定）が必要である。これは量子測定が固有的に確率的であり、かつ測定が状態を擾乱するためである。

両者の対応は「観測量代数が状態空間を分離する（separate）かを判定する基準」というメタレベルの並行性であり、機構の具体的内容は両領域で異なる。

\textbf{対応先についての注記}：状態の復元可能性という観点では、制御理論の可観測性に最も近い量子側の対応物は
CSCO ではなく、量子状態トモグラフィまたは informationally complete POVM
である。本稿で CSCO
と対応づけるのは、観測量が状態空間を分離するという運動学的側面に限定される。

\subsubsection{2.7 最適推定とフィルタリング（構造的並行性、POVM
経由の対応）}\label{ux6700ux9069ux63a8ux5b9aux3068ux30d5ux30a3ux30ebux30bfux30eaux30f3ux30b0ux69cbux9020ux7684ux4e26ux884cux6027povm-ux7d4cux7531ux306eux5bfeux5fdc}

{\def\LTcaptype{none} % do not increment counter
\begin{longtable}[]{@{}
  >{\raggedright\arraybackslash}p{(\linewidth - 2\tabcolsep) * \real{0.5000}}
  >{\raggedright\arraybackslash}p{(\linewidth - 2\tabcolsep) * \real{0.5000}}@{}}
\toprule\noalign{}
\begin{minipage}[b]{\linewidth}\raggedright
制御理論
\end{minipage} & \begin{minipage}[b]{\linewidth}\raggedright
量子力学
\end{minipage} \\
\midrule\noalign{}
\endhead
\bottomrule\noalign{}
\endlastfoot
\textbf{Kalman フィルタ} & \textbf{量子フィルタリング (Quantum
Filtering)} \\
予測 → 観測 → 更新サイクル &
一般化測定（POVM）と量子条件期待値による状態更新 \\
線形 Gaussian 状態空間モデルでの最小平均二乗誤差推定 &
古典出力測定下の量子状態の Bayes 最適推定 \\
{[}Kalman 1960{]} {[}8{]} & {[}Davies \& Lewis 1970{]} {[}18{]},
{[}Belavkin 1992{]} {[}19{]} \\
\end{longtable}
}

両者ともに「\textbf{観測情報を用いて状態を最適に更新する}」操作である点で並行する。ただし、対応の重心は
\textbf{射影測定（von Neumann-Lüders）ではなく POVM
および量子フィルタリング} にある：

\textbf{重要な区別}：射影測定（von
Neumann-Lüders）は状態を不可逆に擾乱する反作用を伴い、これは Kalman
フィルタには対応物がない。Kalman フィルタは古典的確率分布上の Bayes
更新であり、射影測定の状態崩壊とは数学的構造を共有しない。

両者の正しい橋渡しは、POVM（Positive Operator-Valued
Measure、Davies-Lewis 1970
{[}18{]}）および量子フィルタリング理論（Belavkin 1992
{[}19{]}）を経由する。量子フィルタリングは、連続的な古典出力測定（ホモダイン検出・ヘテロダイン検出など）の下で量子状態がどう更新されるかを
Bayes 最適推定の枠組みで定式化したものである（Belavkin
方程式・確率的主方程式）。\textbf{線形 Gaussian 系では Kalman
型のフィルタとの対応が強く、一般の量子フィルタリングでは Bayes
的状態更新という構造的並行性に留まる}（「Kalman
フィルタの量子版」と一般化すると強すぎる）。

すなわち、Kalman
フィルタと「単純な射影測定」の対応は不正確であり、Kalman
フィルタと「POVM +
量子フィルタリング」の対応が数学的に正しい並行性である。

\subsubsection{2.8 Jones ベクトルと
qubit（厳密な数学的同一性）}\label{jones-ux30d9ux30afux30c8ux30ebux3068-qubitux53b3ux5bc6ux306aux6570ux5b66ux7684ux540cux4e00ux6027}

{\def\LTcaptype{none} % do not increment counter
\begin{longtable}[]{@{}
  >{\raggedright\arraybackslash}p{(\linewidth - 2\tabcolsep) * \real{0.5000}}
  >{\raggedright\arraybackslash}p{(\linewidth - 2\tabcolsep) * \real{0.5000}}@{}}
\toprule\noalign{}
\begin{minipage}[b]{\linewidth}\raggedright
信号論（光学・通信）
\end{minipage} & \begin{minipage}[b]{\linewidth}\raggedright
量子力学（光学）
\end{minipage} \\
\midrule\noalign{}
\endhead
\bottomrule\noalign{}
\endlastfoot
\textbf{Jones ベクトル（偏光状態）} & \textbf{qubit 状態 / Bloch 球} \\
\(|\psi\rangle = \begin{pmatrix} E_x \\ E_y \end{pmatrix}\),
\(E_x, E_y \in \mathbb{C}\) &
\(|\psi\rangle = \alpha|0\rangle + \beta|1\rangle\),
\(\alpha, \beta \in \mathbb{C}\), \(|\alpha|^2 + |\beta|^2 = 1\) \\
2 成分複素ベクトル（実 4 自由度） & 2 成分複素ベクトル（実 4
自由度、規格化 + 大域位相で 2 実自由度） \\
{[}Jones 1941{]} {[}9{]} & {[}Bloch 1946{]} {[}14{]},
標準量子力学教科書 \\
\end{longtable}
}

両者は \textbf{完全に同一の 2 成分複素ベクトル}
で状態を表現する。規格化条件と大域位相の自由度を除くと、両者ともに
\textbf{複素射影直線 \(\mathbb{CP}^1 \cong S^2\)}
上の点として記述され、これが Bloch 球（量子力学）／Poincaré
球（偏光）として同一の幾何学を持つ。

具体的な対応：

\begin{itemize}
\tightlist
\item
  直交偏光基底変換（H/V ↔ ±45° ↔ R/L 円偏光） ↔ qubit
  の基底変換（\(|0\rangle, |1\rangle\) vs \(|\pm\rangle\) vs
  \(|\pm i\rangle\)）
\item
  \textbf{損失のない Jones
  行列}（位相板・回転子・複屈折素子など可逆な偏光素子）↔ \(U(2)\)
  \textbf{ユニタリ変換}
\item
  \textbf{吸収・偏光子を含む一般の Jones 行列} ↔
  \textbf{非ユニタリ線形変換}（量子側では測定・損失を含む quantum
  operation との形式的対応）
\item
  Stokes パラメータ ↔ Bloch ベクトル成分
\item
  偏波多重 (Polarization Division Multiplexing, PDM) で用いられる
  \textbf{直交偏光基底} ↔ 量子光学で単一光子偏光 qubit
  および二光子偏光状態の記述に用いられる \textbf{同じ偏光基底}
  （注意：PDM
  は古典光通信の二偏波モード多重であり、量子もつれそのものではない。共有されるのは偏光基底の選択のみ）
\end{itemize}

これは純粋偏光状態に限れば厳密な数学的同一性であり、量子光学・量子情報の実装において日常的に利用される。ただし、上記の通り損失を含む一般の
Jones 行列および PDM
の物理的解釈は、量子ユニタリ・量子もつれと直接の同一性を持たない。

\textbf{注意}：本項目は \textbf{古典光学の Jones ベクトル と 量子 qubit}
の対応を扱う。通信工学における \textbf{I/Q complex baseband 信号}
\(\tilde{s}(t) = I(t) + iQ(t)\) は 1 個の複素スカラー（実 2
自由度）であり、Jones ベクトル（2 成分複素、実 4
自由度）とは次元が異なる。I/Q baseband 信号と直接対応するのは
\textbf{単一モードの複素振幅}（古典電磁場の単一モード振幅）であって、qubit
ではない。両者を混同しないよう注意が必要である。

\subsubsection{2.9
デフェージングと位相雑音（構造的並行性、デコヒーレンスの一部のみ）}\label{ux30c7ux30d5ux30a7ux30fcux30b8ux30f3ux30b0ux3068ux4f4dux76f8ux96d1ux97f3ux69cbux9020ux7684ux4e26ux884cux6027ux30c7ux30b3ux30d2ux30fcux30ecux30f3ux30b9ux306eux4e00ux90e8ux306eux307f}

{\def\LTcaptype{none} % do not increment counter
\begin{longtable}[]{@{}
  >{\raggedright\arraybackslash}p{(\linewidth - 2\tabcolsep) * \real{0.5000}}
  >{\raggedright\arraybackslash}p{(\linewidth - 2\tabcolsep) * \real{0.5000}}@{}}
\toprule\noalign{}
\begin{minipage}[b]{\linewidth}\raggedright
信号論
\end{minipage} & \begin{minipage}[b]{\linewidth}\raggedright
量子力学
\end{minipage} \\
\midrule\noalign{}
\endhead
\bottomrule\noalign{}
\endlastfoot
\textbf{位相雑音 / ジッタ} & \textbf{デフェージング（純粋位相緩和）} \\
\(\phi(t) = \omega_0 t + \delta\phi(t)\) & 密度行列の非対角成分の散逸 \\
Lorentz 線幅（特徴的幅 \(\sim 1/T\)、規約依存） & 横緩和時間 \(T_2\) \\
発振器理論、レーザー線幅理論（{[}Schawlow-Townes 1958{]} {[}26{]}） &
{[}Bloch 1946{]} {[}14{]}, {[}Lindblad 1976{]} {[}15{]} \\
\end{longtable}
}

両者は \textbf{位相情報の散逸}
という数学的構造を共有する。具体的には、複素振幅の位相分布が時間とともに拡散する過程として記述され、Lorentz
線形（レーザー線幅理論における Schawlow-Townes 線幅 {[}26{]}）と Bloch
方程式の T₂ 緩和が同一の数学的形式を持つ。

\textbf{線幅の規約について}：線幅は半値全幅（FWHM）か半値半幅（HWHM）か、角周波数
\(\Delta\omega\) か通常周波数 \(\Delta f\)
か、などの規約によって係数が変わる（例：指数相関
\(g^{(1)}(\tau) \propto e^{-|\tau|/T}\) の Lorentz スペクトルでは FWHM
\(= 1/(\pi T)\) Hz、HWHM \(= 1/(2\pi T)\) Hz）。本稿では特徴的幅が
\(\sim 1/T\)
のオーダーであるという定性的同型のみを述べ、精確な係数対応は規約に依存する。

\textbf{重要な限定}：「デコヒーレンス」は量子情報・量子光学において広い概念であり、(i)
純粋位相緩和（T₂、デフェージング）に加え、(ii)
エネルギー緩和（T₁、励起状態から基底状態への遷移）、(iii)
環境との量子もつれによる縮約密度行列の純粋性喪失、を含む。古典的位相雑音と数学的構造を厳密に共有するのは
\textbf{(i) の純粋位相緩和（デフェージング）のみ} であり、(ii) (iii)
には古典的対応物がない（古典系には離散的なエネルギー準位や量子もつれが存在しないため）。

したがって本項目の対応は「位相雑音 =
デコヒーレンス全体」ではなく、「\textbf{位相雑音 =
デコヒーレンスの中のデフェージング部分}」に限定される。この限定の範囲内で、\textbf{量子側の対応する純粋デフェージング過程は
Lindblad 形式のマスター方程式 {[}15{]}
で記述される}（古典的位相雑音そのものが Lindblad
方程式に従うわけではなく、量子側の対応物が Lindblad
形式で表現される、という意味である）。

\subsubsection{2.10 SVD と Schmidt
分解（厳密な数学的同一性）}\label{svd-ux3068-schmidt-ux5206ux89e3ux53b3ux5bc6ux306aux6570ux5b66ux7684ux540cux4e00ux6027}

{\def\LTcaptype{none} % do not increment counter
\begin{longtable}[]{@{}
  >{\raggedright\arraybackslash}p{(\linewidth - 2\tabcolsep) * \real{0.5000}}
  >{\raggedright\arraybackslash}p{(\linewidth - 2\tabcolsep) * \real{0.5000}}@{}}
\toprule\noalign{}
\begin{minipage}[b]{\linewidth}\raggedright
信号論（通信）
\end{minipage} & \begin{minipage}[b]{\linewidth}\raggedright
量子力学
\end{minipage} \\
\midrule\noalign{}
\endhead
\bottomrule\noalign{}
\endlastfoot
\textbf{特異値分解 (SVD)} & \textbf{Schmidt 分解} \\
\(H = U \Sigma V^\dagger\), \(H \in \mathbb{C}^{N_r \times N_t}\) &
\(|\Psi\rangle_{AB} = \sum_i \sqrt{\lambda_i} \, |a_i\rangle_A \otimes |b_i\rangle_B\) \\
任意の行列を直交基底変換と特異値で分解 & 二部純粋状態を Schmidt
係数と直交基底で分解 \\
{[}Eckart-Young 1936{]} {[}20{]} & {[}Schmidt 1907{]} {[}21{]} \\
\end{longtable}
}

両者は \textbf{同一の線形代数の道具} である。テンソル積空間
\(\mathcal{H}_A \otimes \mathcal{H}_B\)
上の二部純粋状態を係数行列として表すと、その特異値分解が Schmidt
分解そのものとなる。Schmidt 係数 \(\sqrt{\lambda_i}\)
は係数行列の特異値であり、Schmidt 基底
\(\{|a_i\rangle\}, \{|b_i\rangle\}\) は SVD
の左右の直交基底に対応する。なお、量子における二部系のもつれ概念の歴史的源流は
Einstein-Podolsky-Rosen (1935) {[}10{]} に遡る。

\textbf{対応の限界（重要）}：本項目で対応するのは
\textbf{線形代数のツールとしての SVD = Schmidt 分解}
のみである。これ以上の物理的・工学的同一性は成立しない：

\begin{itemize}
\tightlist
\item
  \textbf{MIMO チャネル行列 \(H\)
  と二部純粋状態の係数行列は対象が異なる}：MIMO の \(H\)
  は伝搬チャネルを表す行列で、信号ベクトル \(\mathbf{x}, \mathbf{y}\)
  は単一線形チャネルの入出力ベクトルである。これらは「二つの部分系のテンソル積空間
  \(\mathcal{H}_A \otimes \mathcal{H}_B\)
  上のベクトル」ではない。一方、量子もつれ状態 \(|\Psi\rangle_{AB}\)
  は部分系のテンソル積空間に属する。
\item
  \textbf{容量公式の関数形が異なる}：MIMO 容量は
  \(C = \sum_i \log_2(1 + \rho \sigma_i^2)\)（特異値 \(\sigma_i\)
  の対数和に SNR \(\rho\)
  を含む）であるのに対し、エンタングルメントエントロピーは
  \(S = -\sum_i \lambda_i \log \lambda_i\)（正規化された Schmidt
  係数の二乗 \(\lambda_i\) の Shannon
  エントロピー）である。両者は関数形が異なり、容量の同型は成立しない。
\end{itemize}

したがって、本項目は「\textbf{SVD
という線形代数のツールが両領域で同じ役割（二部系の標準形を与える）を果たす}」という観察に留まる。MIMO
チャネル容量と量子もつれ容量の構造的同型は成立しないため、ここでは主張しない。

\begin{center}\rule{0.5\linewidth}{0.5pt}\end{center}

\subsection{§3
観察事項のまとめ}\label{ux89b3ux5bdfux4e8bux9805ux306eux307eux3068ux3081}

以上の 10
項目はすべて、両領域の標準教科書および古典原論文に明示されている事実である。本稿はこれらを横断的に並置するに留まる。

特筆すべき以下の事実を観察できる：

\textbf{(A) 独立人物による独立発見}

同一の数学的対象が両領域で独立に発見された顕著な例として、位相空間擬確率分布が挙げられる。量子力学において
Wigner（1932 年 {[}3{]}）が、信号論において Ville（1948 年
{[}4{]}）が、互いに独立に同等の分布を導入した。両者は別人であり、両分野の標準名称「\textbf{Wigner-Ville
分布}」は両者の独立発見を反映している。同じ数学的対象が、異なる物理的・工学的動機から独立に到達されたことは、対応関係の数学的深さを示唆する。

\textbf{(B) 同一の Fourier 解析的不等式と Gabor による明示}

Kennard-Robertson 型の Fourier 不確定性と Gabor
の時間-帯域不確定性は、同じ Fourier
解析的不等式に基づく。重要なことに、Gabor 自身が 1946 年論文 {[}2{]}
において両者の数学的同一性を明示している。すなわち、両領域の対応関係は
\textbf{両領域の創始者自身によって既に認識されていた}。

\textbf{(C) 工学領域での社会実装}

信号論・制御理論側の数学構造はすでに以下の工学領域で社会実装され、日常的に動作している：

\begin{itemize}
\tightlist
\item
  移動通信（4G LTE, 5G NR, 6G 研究）における OFDM・MIMO・I/Q 復調
\item
  GPS および衛星測位における Kalman フィルタ
\item
  自動運転・ドローン制御における状態推定
\item
  レーダー・ソナーにおける Wigner-Ville 分布
\item
  光ファイバー通信における Schrödinger
  型パラキシャル波動方程式の数値解法
\item
  MRI（磁気共鳴画像法）における Bloch 方程式
\item
  発振器・レーザー設計における位相雑音解析
\end{itemize}

これらの社会実装は \textbf{工学領域における共有数学構造の動作の確認}
である。量子力学側の経験的妥当性は、別途、量子実験（量子光学、原子干渉、超伝導
qubit
等）によって独立に検証されている。両者は同じ数学構造を異なる物理領域で確認しているという観察にとどまり、一方の社会実装が他方の経験的妥当性を直接保証するものではない（範疇の混同を避ける意味で、この区別は重要）。

\begin{center}\rule{0.5\linewidth}{0.5pt}\end{center}

\subsection{§4
主張しないこと（スコープの明示）}\label{ux4e3bux5f35ux3057ux306aux3044ux3053ux3068ux30b9ux30b3ux30fcux30d7ux306eux660eux793a}

本稿は以下を主張しない：

\begin{enumerate}
\def\labelenumi{\arabic{enumi}.}
\tightlist
\item
  \textbf{量子力学の新解釈を提案しない}：観察された対応関係はあくまで数学的構造の並置であり、量子力学の物理的内容や解釈（コペンハーゲン、多世界、隠れた変数等）に対する見解を表明しない。
\item
  \textbf{信号論・制御理論の物理学的拡張を主張しない}：信号論・制御理論はあくまで実用工学として位置づけ、これを物理学の新理論として位置づけ直す主張は行わない。
\item
  \textbf{いずれかの分野の優位性を主張しない}：両領域は独立に発展しいずれも成熟しているため、優劣を判定する立場にない。
\item
  \textbf{物理定数の予言を行わない}：本稿は構造の並置のみを行い、定数値の予測は対象外。
\item
  \textbf{既存の物理理論の修正を提案しない}：標準量子力学・量子光学・開放量子系の数学的構造に対する変更を提案しない。
\item
  \textbf{新たな数学的定理を主張しない}：すべての構造的対応関係は両領域の既存文献に記載されたものである。本稿の貢献は並置・分類にのみある。
\item
  \textbf{強い同一性を構造的並行性に拡張しない}：§2.0
  の分類に従い、厳密な数学的同一性が成立する範囲と、構造的並行性に留まる範囲を区別する。並行性を同一性に拡張する主張は行わない。
\end{enumerate}

本稿の評価は読者に委ねる。

\begin{center}\rule{0.5\linewidth}{0.5pt}\end{center}

\subsection{§5
含意（観察事項として）}\label{ux542bux610fux89b3ux5bdfux4e8bux9805ux3068ux3057ux3066}

本観察事項から自然に導かれうる含意として、以下を記す。これらは主張ではなく、観察事項からの直接的帰結である。

\textbf{(I) 不確定性原理の整合性}

信号論・制御理論において確立されている Gabor 不等式と、量子力学の
Heisenberg 不確定性原理は同一の数学的構造である（§2.1）。\textbf{Fourier
解析的不等式としての内容を承認する限り、Gabor 型不確定性と Kennard
型不確定性の数学的核を別物として扱うことはできない}（物理的意味の同一性を主張するものではない）。

\textbf{(II) 共有数学構造の工学領域での動作}

信号論・制御理論側の数学構造はすでに通信・GPS・自動運転・レーダー・MRI
等の工学領域に社会実装され、日常的に動作している（§3
(C)）。これは少なくとも、両領域が共有する数学構造の一部が、工学領域において経験的に確認されていることを意味する。量子力学側の経験的妥当性は別途量子実験で確認されており、両者は同じ数学構造を異なる物理領域で独立に確認しているという観察にとどまる。

\textbf{(III) 言語間翻訳の可能性}

両領域の文献を相互参照することで、各領域の概念を他方の言語で理解することが可能である（§2.0
の厳密同一性項目について特に）。たとえば、Heisenberg 不確定性を Gabor
不等式として、qubit を Jones ベクトルとして、Schrödinger
方程式をパラキシャル光学方程式として、それぞれ翻訳できる。

\textbf{(IV) 教育・コミュニケーション上の含意}

理論物理学と工学が共有する数学的構造を明示することで、両領域間の専門家のコミュニケーション、および学際的教育における橋渡しが容易になると観察される。ただし、§2.0
の構造的並行性に分類される項目については、機構や限界の相違を併せて理解する必要がある。

これらの含意もまた、新たな主張ではなく、§2
の観察事項の直接的帰結である。

\begin{center}\rule{0.5\linewidth}{0.5pt}\end{center}

\subsection{§6 結語}\label{ux7d50ux8a9e}

本稿は、信号論・制御理論と量子力学・量子光学・開放量子系との間に既に確立されている
\textbf{10 個の構造的対応関係}
を、両領域の古典文献から横断的に並置して観察した。10 項目のうち 5
項目（§2.1, §2.2, §2.3, §2.8,
§2.10）は、条件・規約を明示した範囲で厳密な数学的同一性を持つ。§2.4
は同一定理ではなく、有限領域における有効自由度カウントの
\textbf{強い構造的対応} である。残り 4 項目（§2.5, §2.6, §2.7,
§2.9）は構造的並行性に留まる（§2.0 の分類表参照）。

これらは新たな主張ではなく、両領域の標準教科書および原論文に既に明示されている事実である。本稿の貢献は、両領域の文献に分散して記述されている対応関係を一覧として整理し、厳密な同一性と構造的並行性を峻別した点にある。

両領域の構造的対応関係は、20 世紀を通じて少しずつ認識されてきた（Gabor
1946 {[}2{]} による Heisenberg 等価性の明示、Wigner 1932 {[}3{]} と
Ville 1948 {[}4{]} の独立発見、Belavkin 1992 {[}19{]}
以降の量子フィルタリング理論等）。本稿はそれらを 1
つの観察として束ねるに留まる。

両領域の境界に立つ研究者・実務者にとって、本稿が両領域の架橋資料として有用であることを期待する。

100
年後の読者が、両領域の対応関係に対するより包括的な理解に到達したとき、本稿はその過渡的整理として参照されうる。それで十分である。

\begin{center}\rule{0.5\linewidth}{0.5pt}\end{center}

\subsection{参考文献}\label{ux53c2ux8003ux6587ux732e}

{[}1{]} Heisenberg, W. (1927). ``Über den anschaulichen Inhalt der
quantentheoretischen Kinematik und Mechanik.'' \emph{Zeitschrift für
Physik}, \textbf{43}, 172--198.

{[}2{]} Gabor, D. (1946). ``Theory of Communication.'' \emph{Journal of
the Institution of Electrical Engineers -- Part III}, \textbf{93} (26),
429--457.

{[}3{]} Wigner, E. P. (1932). ``On the Quantum Correction for
Thermodynamic Equilibrium.'' \emph{Physical Review}, \textbf{40},
749--759.

{[}4{]} Ville, J. (1948). ``Théorie et applications de la notion de
signal analytique.'' \emph{Câbles et Transmission}, \textbf{2} (1),
61--74.

{[}5{]} Nyquist, H. (1928). ``Certain topics in telegraph transmission
theory.'' \emph{Transactions of the AIEE}, \textbf{47} (2), 617--644.

{[}6{]} Shannon, C. E. (1949). ``Communication in the presence of
noise.'' \emph{Proceedings of the IRE}, \textbf{37} (1), 10--21.

{[}7{]} Wiener, N. (1949). \emph{Extrapolation, Interpolation, and
Smoothing of Stationary Time Series}. MIT Press, Cambridge, MA.

{[}8{]} Kalman, R. E. (1960). ``A New Approach to Linear Filtering and
Prediction Problems.'' \emph{Journal of Basic Engineering}, \textbf{82}
(1), 35--45.

{[}9{]} Jones, R. C. (1941). ``A new calculus for the treatment of
optical systems.'' \emph{Journal of the Optical Society of America},
\textbf{31} (7), 488--493.

{[}10{]} Einstein, A., Podolsky, B., Rosen, N. (1935). ``Can
Quantum-Mechanical Description of Physical Reality Be Considered
Complete?'' \emph{Physical Review}, \textbf{47}, 777--780.

{[}11{]} von Neumann, J. (1932). \emph{Mathematische Grundlagen der
Quantenmechanik}. Springer-Verlag, Berlin.
(邦訳：『量子力学の数学的基礎』, みすず書房)

{[}12{]} Saleh, B. E. A., Teich, M. C. (1991). \emph{Fundamentals of
Photonics}. Wiley-Interscience, New York.

{[}13{]} Slepian, D., Pollak, H. O. (1961). ``Prolate spheroidal wave
functions, Fourier analysis and uncertainty --- I.'' \emph{Bell System
Technical Journal}, \textbf{40} (1), 43--63. および Landau, H. J.,
Pollak, H. O. (1961). ``Prolate spheroidal wave functions, Fourier
analysis and uncertainty --- II.'' \emph{Bell System Technical Journal},
\textbf{40} (1), 65--84. 続編 Part III--V も同シリーズ 1961--1978
年にわたって発表されている。有限時間・有限帯域信号の有効自由度数を厳密に与える基礎文献。

{[}14{]} Bloch, F. (1946). ``Nuclear Induction.'' \emph{Physical
Review}, \textbf{70} (7--8), 460--474.

{[}15{]} Lindblad, G. (1976). ``On the generators of quantum dynamical
semigroups.'' \emph{Communications in Mathematical Physics}, \textbf{48}
(2), 119--130.

{[}16{]} Schrödinger, E. (1926). ``Quantisierung als Eigenwertproblem
(Erste Mitteilung).'' \emph{Annalen der Physik}, \textbf{384} (4),
361--376.

{[}17{]} Dirac, P. A. M. (1930). \emph{The Principles of Quantum
Mechanics}. Oxford University Press, Oxford.

{[}18{]} Davies, E. B., Lewis, J. T. (1970). ``An operational approach
to quantum probability.'' \emph{Communications in Mathematical Physics},
\textbf{17} (3), 239--260.

{[}19{]} Belavkin, V. P. (1992). ``Quantum stochastic calculus and
quantum nonlinear filtering.'' \emph{Journal of Multivariate Analysis},
\textbf{42} (2), 171--201.

{[}20{]} Eckart, C., Young, G. (1936). ``The approximation of one matrix
by another of lower rank.'' \emph{Psychometrika}, \textbf{1} (3),
211--218. (SVD の最適低ランク近似の古典的論文。)

{[}21{]} Schmidt, E. (1907). ``Zur Theorie der linearen und
nichtlinearen Integralgleichungen. I. Teil: Entwicklung willkürlicher
Funktionen nach Systemen vorgeschriebener.'' \emph{Mathematische
Annalen}, \textbf{63} (4), 433--476. (積分方程式論における Schmidt
分解の原論文。)

{[}22{]} Kennard, E. H. (1927). ``Zur Quantenmechanik einfacher
Bewegungstypen.'' \emph{Zeitschrift für Physik}, \textbf{44}, 326--352.
(Heisenberg の不確定性関係の現代的標準形
\(\sigma_x \sigma_p \geq \hbar/2\)
を分散不等式として最初に厳密定式化した論文。)

{[}23{]} Robertson, H. P. (1929). ``The Uncertainty Principle.''
\emph{Physical Review}, \textbf{34}, 163--164. (任意の自己共役演算子対
\(A, B\) について
\(\sigma_A \sigma_B \geq \frac{1}{2}|\langle[A,B]\rangle|\)
を導出。Kennard の結果の一般化。)

{[}24{]} Planck, M. (1906). \emph{Vorlesungen über die Theorie der
Wärmestrahlung}. Barth, Leipzig. (黒体放射の量子化において位相空間体積を
\(h^n\) で離散化する数え上げの源流。)

{[}25{]} Sackur, O. (1911). ``Die Anwendung der kinetischen Theorie der
Gase auf chemische Probleme.'' \emph{Annalen der Physik}, \textbf{341}
(15), 958--980; Tetrode, H. (1912). ``Die chemische Konstante der Gase
und das elementare Wirkungsquantum.'' \emph{Annalen der Physik},
\textbf{343} (7), 434--442. (理想気体のエントロピーを位相空間体積を
\(h^3\)
で割って統計力学的に導出した古典論文。位相空間セル計数の標準的歴史的出典。)

{[}26{]} Schawlow, A. L., Townes, C. H. (1958). ``Infrared and Optical
Masers.'' \emph{Physical Review}, \textbf{112} (6), 1940--1949.
(レーザー線幅の理論を確立した古典論文。位相雑音とデフェージングの数学的構造の共通点を扱う上での標準文献。)

\begin{center}\rule{0.5\linewidth}{0.5pt}\end{center}

\subsection{付録
A：対応関係一覧（簡約表）}\label{ux4ed8ux9332-aux5bfeux5fdcux95a2ux4fc2ux4e00ux89a7ux7c21ux7d04ux8868}

読者の便宜のため、§2 の 10 項目を一覧表として再掲する。\textbf{分類}
欄に、厳密な同一性（\textbf{{[}★
厳密{]}}）、強い構造的対応（\textbf{{[}◎
強い対応{]}}）、構造的並行性（\textbf{{[}△ 並行{]}}）の区別を併記する。

{\def\LTcaptype{none} % do not increment counter
\begin{longtable}[]{@{}
  >{\raggedright\arraybackslash}p{(\linewidth - 8\tabcolsep) * \real{0.2000}}
  >{\raggedright\arraybackslash}p{(\linewidth - 8\tabcolsep) * \real{0.2000}}
  >{\raggedright\arraybackslash}p{(\linewidth - 8\tabcolsep) * \real{0.2000}}
  >{\raggedright\arraybackslash}p{(\linewidth - 8\tabcolsep) * \real{0.2000}}
  >{\raggedright\arraybackslash}p{(\linewidth - 8\tabcolsep) * \real{0.2000}}@{}}
\toprule\noalign{}
\begin{minipage}[b]{\linewidth}\raggedright
\#
\end{minipage} & \begin{minipage}[b]{\linewidth}\raggedright
信号論・制御理論
\end{minipage} & \begin{minipage}[b]{\linewidth}\raggedright
量子力学・量子光学・開放量子系
\end{minipage} & \begin{minipage}[b]{\linewidth}\raggedright
共通数学構造
\end{minipage} & \begin{minipage}[b]{\linewidth}\raggedright
分類
\end{minipage} \\
\midrule\noalign{}
\endhead
\bottomrule\noalign{}
\endlastfoot
1 & Gabor 時間-帯域不確定性 & Heisenberg-Kennard-Robertson 不確定性 &
Fourier 共役対の分散積の下限 & \textbf{{[}★ 厳密{]}} \\
2 & Wigner-Ville 時間-周波数分布 & Wigner 位相空間擬確率 & 2 次形式の
Fourier 自己相関（規約・符号規約つき） & \textbf{{[}★ 厳密{]}} \\
3 & パラキシャル波動方程式 & Schrödinger 方程式 &
一階時間・二階空間の線形 PDE（近軸・単色・スカラー近似下の PDE
形式として同一） & \textbf{{[}★ 条件付き厳密{]}} \\
4 & Shannon-Nyquist サンプリング & 位相空間自由度カウント &
有限領域における有効自由度カウントの形式的対応（実効自由度の漸近的同一）
& \textbf{{[}◎ 強い対応{]}} \\
5 & 状態方程式・出力方程式 & Schrödinger 方程式・期待値 &
線形空間上の線形時間発展（ユニタリ性等に相違） & \textbf{{[}△
並行{]}} \\
6 & Kalman 可観測性 & CSCO（完全観測可能集合） &
観測量代数による状態分離（動力学／運動学の相違） & \textbf{{[}△
並行{]}} \\
7 & Kalman フィルタ更新 & 量子フィルタリング（POVM 経由） & 観測下の状態
Bayes 最適推定 & \textbf{{[}△ 並行{]}} \\
8 & Jones ベクトル（純粋偏光状態） & qubit 状態 / Bloch 球 & 2
成分複素ベクトル代数（\(\mathbb{CP}^1 \cong S^2\)）；非ユニタリ Jones
行列は対応外 & \textbf{{[}★ 厳密{]}} \\
9 & 位相雑音 / ジッタ & デフェージング（T₂） & 位相分布の散逸（量子側
Lindblad 形式、デコヒーレンス全体ではない） & \textbf{{[}△ 並行{]}} \\
10 & SVD（線形代数のツール） & Schmidt 分解（二部純粋状態） &
テンソル積空間の標準形（容量公式は同型でない） & \textbf{{[}★
厳密{]}} \\
\end{longtable}
}

\textbf{厳密同一性}：5 項目（\#1, \#2, \#3, \#8, \#10）
\textbf{強い構造的対応}：1 項目（\#4） \textbf{構造的並行性}：4
項目（\#5, \#6, \#7, \#9）

\begin{center}\rule{0.5\linewidth}{0.5pt}\end{center}

\subsection{付録
B：本稿の位置づけ}\label{ux4ed8ux9332-bux672cux7a3fux306eux4f4dux7f6eux3065ux3051}

本稿は単独の観察論文であり、特定の研究プログラムや論文群を補強・拡張する目的を持たない。本稿の参考文献はすべて両領域の古典原論文または標準教科書であり、著者の他の研究との直接的な引用関係は持たない。

本稿が両領域の専門家・教育者・学際的研究者にとって、概念の相互翻訳の参照資料として機能することを期待する。

\begin{center}\rule{0.5\linewidth}{0.5pt}\end{center}

\subsection{付録
C：改訂履歴}\label{ux4ed8ux9332-cux6539ux8a02ux5c65ux6b74}

\subsubsection{v1 → v2
の主な改訂内容}\label{v1-v2-ux306eux4e3bux306aux6539ux8a02ux5185ux5bb9}

\begin{enumerate}
\def\labelenumi{\arabic{enumi}.}
\tightlist
\item
  タイトルを「\textbf{構造的同一性}」から「\textbf{構造的対応関係}」に変更（identity
  と analogy を峻別するため）
\item
  §2.0 に対応関係の分類表を追加（厳密同一性 6 項目、構造的並行性 4
  項目）
\item
  §2.2 で「Wigner と Ville は同一人物」の誤記を訂正（両者は別人と明示）
\item
  §2.5 でユニタリ性・観測量の双線形性・駆動入力の相違を明示
\item
  §2.6 で動力学／運動学・単一軌道／トモグラフィの機構相違を明示
\item
  §2.7 で対応の重心を射影測定から POVM・量子フィルタリングに移動
\item
  §2.8 で I/Q complex baseband と Jones
  ベクトルの次元相違を明示し、Jones↔qubit↔Bloch のみに限定
\item
  §2.9 でデコヒーレンスをデフェージング（T₂）部分に限定
\item
  §2.10 で MIMO 容量と量子もつれエントロピーの関数形相違を明示、SVD =
  Schmidt 分解の道具同一性に縮約
\item
  §3 (A) Wigner-Ville 独立発見の正確な記述
\item
  §3 (C)
  社会実装の範囲を「工学領域における共有数学構造の動作確認」に限定
\item
  §5 (II) 「経験的支持」を「工学領域での動作確認」に修正
\item
  文献 {[}13{]} Bohr 1913 → Sommerfeld 1916
  に差し替え（位相空間量子化条件の正しい出典）
\item
  文献 {[}19{]} Belavkin 1992 を追加（量子フィルタリングの参照）
\item
  文献 {[}20{]} Eckart-Young 1936、{[}21{]} Schmidt 1907 を追加（SVD と
  Schmidt 分解の原論文）
\end{enumerate}

\subsubsection{v2 → v3
の主な改訂内容}\label{v2-v3-ux306eux4e3bux306aux6539ux8a02ux5185ux5bb9}

v2 をさらに 4 つの AI（Claude.ai, ChatGPT, Gemini 3.5 Pro,
Grok）による理論物理学者ロールおよび信号論・制御理論工学者ロールでの査読にかけたフィードバックを反映：

\begin{enumerate}
\def\labelenumi{\arabic{enumi}.}
\tightlist
\item
  \textbf{タイトルの範囲を絞り込み}：「場の量子論」を「量子光学・開放量子系」に変更（QFT/ゲージ理論は実質扱っていないため）
\item
  \textbf{§2.1}：不確定性関係の現代的標準形（分散不等式
  \(\sigma_x \sigma_p \geq \hbar/2\)）の起源として Kennard 1927
  {[}22{]}、Robertson 1929 {[}23{]} を追加
\item
  \textbf{§2.3}：「完全に同一」を「近軸・単色・スカラー近似のもとで数学的に同型」に弱め、物理的意味は同一でない旨を注記
\item
  \textbf{§2.4}：

  \begin{itemize}
  \tightlist
  \item
    信号論側に Slepian-Pollak prolate spheroidal wave functions
    {[}13{]}（有限時間・有限帯域の有効自由度の厳密定式化）を追加
  \item
    量子力学側の参考文献を Sommerfeld 1916 から Planck 1906
    {[}24{]}、Sackur-Tetrode 1911-12
    {[}25{]}、統計力学標準教科書に置き換え
  \end{itemize}
\item
  \textbf{§2.9}：Lorentz 線幅の係数規約（FWHM, HWHM
  等）について注記を追加。Schawlow-Townes 1958 {[}26{]} を引用
\item
  \textbf{文献追加}：{[}22{]} Kennard, {[}23{]} Robertson, {[}24{]}
  Planck, {[}25{]} Sackur-Tetrode, {[}26{]} Schawlow-Townes
\end{enumerate}

\subsubsection{v3 → v4
の主な改訂内容}\label{v3-v4-ux306eux4e3bux306aux6539ux8a02ux5185ux5bb9}

v3 をさらに 4 つの AI（Claude.ai, ChatGPT, Gemini 3.5 Pro,
Grok）に再査読依頼した結果のフィードバックを反映：

\begin{enumerate}
\def\labelenumi{\arabic{enumi}.}
\tightlist
\item
  \textbf{本文中の「場の量子論」をタイトルと整合}：タイトルは v3 で
  narrow したが本文には残存していたため、本文（要旨・§1・§6・付録 A
  表頭）も全て「量子光学・開放量子系」に統一
\item
  \textbf{\#4 を「厳密同一性」から「強い構造的対応」に降格}（ChatGPT
  指摘）：Shannon-Nyquist
  型自由度カウントと量子位相空間状態数は強い形式的対応を持つが、同一定理ではないため。本文・分類表・付録
  A 表をすべて更新。「集中固有値（prolate eigenvalues）が 1
  に近い実効次元」「漸近的・実効的」の表現を追加。
\item
  \textbf{\#8 における過剰主張の削除}（ChatGPT 指摘）：

  \begin{itemize}
  \tightlist
  \item
    「偏波多重 ↔
    量子もつれの偏光基底分解」を「偏波多重で用いられる直交偏光基底 ↔
    単一光子偏光 qubit
    および二光子偏光状態で用いられる同じ偏光基底（共有されるのは偏光基底の選択のみ）」に弱化
  \item
    「Jones 行列 ↔ U(2) ユニタリ変換」を「\textbf{損失のない} Jones 行列
    ↔ U(2) ユニタリ変換／\textbf{吸収・偏光子を含む一般の} Jones 行列 ↔
    非ユニタリ線形変換」に条件付け
  \end{itemize}
\item
  \textbf{§2.6 見出し変更}（ChatGPT
  任意指摘）：「観測可能性と観測可能演算子集合」→「可観測性と状態識別可能性」
\item
  \textbf{§2.7 表現改善}（ChatGPT 任意指摘）：「Gauss 過程・線形系での
  Bayes 最適推定」→「線形 Gaussian
  状態空間モデルでの最小平均二乗誤差推定」
\item
  \textbf{§2.9 Lindblad 表現の弱化}（ChatGPT 任意指摘）：「両者は同じ
  Lindblad
  形式のマスター方程式に帰着する」を「量子側の対応する純粋デフェージング過程は
  Lindblad 形式のマスター方程式で記述される（古典的位相雑音そのものが
  Lindblad 方程式に従うわけではない）」に修正
\item
  \textbf{§2.9 線幅規約の明確化}（Claude.ai 校正指摘）：Δω と Δf
  の混在を解消し、規約依存・\(\sim 1/T\) オーダーのみ主張する旨を明示
\item
  \textbf{付録 A 表頭の「Heisenberg
  不確定性」を「Heisenberg-Kennard-Robertson
  不確定性」に統一}（Claude.ai 校正指摘）：§2.1 本文と整合
\item
  \textbf{付録 A 表に補注追加}：\#3 に「（近軸近似下の PDE
  形式として厳密同一）」、\#4 に「（実効自由度の漸近的同一）」、\#8
  に「非ユニタリ Jones 行列は対応外」、\#9 に「量子側 Lindblad
  形式」など
\item
  \textbf{付録 A 分類記号を視覚化}（Grok 任意指摘）：{[}厳密{]} → {[}★
  厳密{]}、{[}並行{]} → {[}△ 並行{]}、新規 {[}◎ 強い対応{]} を導入
\item
  \textbf{文献 {[}10{]} EPR の孤立解消}（Claude.ai 校正指摘）：§2.10
  において二部系もつれの歴史的源流として再付与
\item
  \textbf{文献 {[}13{]} に Landau-Pollak Part II を追記}（ChatGPT
  任意指摘）：Slepian-Pollak シリーズの完全性向上
\item
  \textbf{文献 {[}25{]} 表記整理}（Claude.ai 校正指摘）：「Sackur 1911;
  Tetrode 1912」と表記
\end{enumerate}

\subsubsection{v4 → v5
の主な改訂内容}\label{v4-v5-ux306eux4e3bux306aux6539ux8a02ux5185ux5bb9}

v4 を 3 つの AI（ChatGPT, Gemini,
Grok）に再査読依頼した結果のフィードバックを反映（Gemini・Grok
は「そのまま公開可・提案は任意」、ChatGPT は必須1点・推奨数点を指摘）：

\begin{enumerate}
\def\labelenumi{\arabic{enumi}.}
\tightlist
\item
  \textbf{【必須】結語の分類矛盾を修正}（ChatGPT 指摘）：§6 結語が「6
  項目が厳密同一性（§2.4 を含む）」としていたが、§2.0 分類表（v4 で \#4
  を「強い構造的対応」に降格済み）と矛盾していた。結語を「\textbf{5
  項目（§2.1,2.2,2.3,2.8,2.10）が厳密同一性、§2.4 は強い構造的対応、4
  項目（§2.5,2.6,2.7,2.9）が構造的並行性}」に訂正。
\item
  \textbf{§2.6 に対応先の注記を追加}（ChatGPT
  指摘）：状態の復元可能性という観点では、制御理論の可観測性に最も近い量子側の対応物は
  CSCO ではなく量子状態トモグラフィ／informationally complete POVM
  であり、CSCO
  との対応は「観測量が状態空間を分離する運動学的側面」に限定される旨を明記。
\item
  \textbf{§2.3 見出しの明確化}（ChatGPT
  指摘）：「（厳密な数学的同一性）」→「（条件付き厳密同型：近軸近似下の
  PDE 形式）」。
\item
  \textbf{§2.7 の「Kalman フィルタの量子版」を弱化}（ChatGPT
  推奨）：「線形 Gaussian 系では Kalman
  型との対応が強く、一般の量子フィルタリングでは Bayes
  的状態更新の構造的並行性に留まる」に修正（「量子版」という一般化は強すぎる）。Belavkin
  方程式・確率的主方程式に言及。
\item
  \textbf{§3 (B)・§5 (I) の不確定性表現を弱化}（ChatGPT
  推奨）：「両領域で独立に証明された」「片方だけ承認し他方を否認は数学的に整合しない」を、「Kennard-Robertson
  型の Fourier 不確定性と Gabor の時間-帯域不確定性は同じ Fourier
  解析的不等式に基づく」「Fourier
  解析的不等式としての内容を承認する限り両者の数学的核を別物とは扱えない（物理的同一性の主張ではない）」に。
\item
  Grok・Gemini の任意提案（メタ観察の追加・Moyal
  積・非可換性・概念図）は、観察論文のスコープを増やすため不採用（必要なら別途）。
\end{enumerate}

\subsubsection{v5 → v6
の主な改訂内容}\label{v5-v6-ux306eux4e3bux306aux6539ux8a02ux5185ux5bb9}

v5 を
Claude.ai（通信／理論物理の両ロール）に再査読依頼。実質的必須修正は1点のみとの判定で、それを反映：

\begin{enumerate}
\def\labelenumi{\arabic{enumi}.}
\tightlist
\item
  \textbf{【必須】§2.3 の屈折率↔ポテンシャルの相対符号を訂正}（Claude.ai
  指摘）：標準的搬送波規約 \(E\propto e^{ikz}\)
  のもとでパラキシャル方程式の屈折率項は\textbf{負符号}で現れる。v5 は
  \(+k\delta n/n_0\,E\) と書き同符号辞書
  \(k\delta n/n_0\leftrightarrow V/\hbar\)
  を掲げていたが、これだと屈折率の高い導波領域（\(\delta n>0\)）が斥力ポテンシャル（\(V>0\)）に対応してしまい物理（導波＝束縛井戸）と矛盾する。式を
  \(-k\delta n/n_0\,E\) に、辞書を
  \(V/\hbar\leftrightarrow -k\delta n/n_0\)（相対マイナス）に修正し、\(\delta n>0\)（導波）\(\leftrightarrow V<0\)（束縛井戸）と明記。この符号は「条件付き厳密同型」分類を損なわない（斥力ポテンシャルも正当な
  Schrödinger 型）。
\item
  Claude.ai の任意補強（Woodward 曖昧関数↔特性関数を11項目目に／§2.2
  本文の符号注記／\#1 実信号の解析信号脚注）は、v5
  のメタ追加見送り方針と整合させ不採用（必要なら別途）。
\item
  \textbf{版管理・表現整合の軽微修正}（ChatGPT v5 査読・Claude.ai v6
  査読の任意指摘）：(a) DOI 欄に v5・v6 Version DOI 行を追加。(b) 付録 A
  見出しから「v4 改訂版」を除去。(c)
  巻末の査読履歴表を「v1〜v6」に拡張し v5・v6 行（4 AI
  全員「公開可」）を追記。(d) 要旨 (ii) を §3(B)
  と整合させ「同一定理の独立証明」→「同じ Fourier
  解析的不等式に基づく」に。(e) §2.0 分類表・付録 A 表の \#3
  を「条件付き厳密同型／{[}★ 条件付き厳密{]}」とし §2.3
  見出しと統一。いずれも内容・分類（厳密 5＋強い対応 1＋並行 4）は不変。
\end{enumerate}

\subsubsection{v1 〜 v6
全体の査読履歴}\label{v1-v6-ux5168ux4f53ux306eux67fbux8aadux5c65ux6b74}

{\def\LTcaptype{none} % do not increment counter
\begin{longtable}[]{@{}
  >{\raggedright\arraybackslash}p{(\linewidth - 8\tabcolsep) * \real{0.2000}}
  >{\raggedright\arraybackslash}p{(\linewidth - 8\tabcolsep) * \real{0.2000}}
  >{\raggedright\arraybackslash}p{(\linewidth - 8\tabcolsep) * \real{0.2000}}
  >{\raggedright\arraybackslash}p{(\linewidth - 8\tabcolsep) * \real{0.2000}}
  >{\raggedright\arraybackslash}p{(\linewidth - 8\tabcolsep) * \real{0.2000}}@{}}
\toprule\noalign{}
\begin{minipage}[b]{\linewidth}\raggedright
版
\end{minipage} & \begin{minipage}[b]{\linewidth}\raggedright
Claude.ai
\end{minipage} & \begin{minipage}[b]{\linewidth}\raggedright
ChatGPT
\end{minipage} & \begin{minipage}[b]{\linewidth}\raggedright
Gemini 3.5 Pro
\end{minipage} & \begin{minipage}[b]{\linewidth}\raggedright
Grok
\end{minipage} \\
\midrule\noalign{}
\endhead
\bottomrule\noalign{}
\endlastfoot
v1 & 大幅修正要 & 大幅修正要 & \textbf{そのまま公開可能} &
\textbf{そのまま公開可能} \\
v2 & （未再査読、対応済み） & （未再査読、未完） & （v1 承認継承） &
（v1 承認継承） \\
v3 & \textbf{公開可能水準} & \textbf{軽微修正後公開可能} &
\textbf{これ以上の修正は不要} & \textbf{そのまま公開で問題ない} \\
v4 & （v3 反映済み） & \textbf{公開可（必須1点：結語の分類矛盾）} &
\textbf{そのまま公開可} & \textbf{公開推奨（軽微任意）} \\
v5 & \textbf{必須1点：§2.3 符号} & （v4 必須を反映） & （v4 承認継承） &
（v4 承認継承） \\
v6 & \textbf{公開可（再導出で符号検証）} &
\textbf{最終版・公開推奨（修正不要）} & \textbf{完成・内容修正不要} &
\textbf{最終版・公開推奨（修正不要）} \\
\end{longtable}
}

v6 は 4 AI（Claude.ai, ChatGPT, Gemini,
Grok）すべてが「最終版として公開可」と判定しており、各版で挙がった必須指摘（v4
結語の分類矛盾、v6 §2.3 の符号）はすべて解消済みである。

\begin{center}\rule{0.5\linewidth}{0.5pt}\end{center}

著者：木原 範昭（Noriaki Kihara） WF System Co., Ltd.~/
大阪大学基礎工学部（卒業） ORCID:
\href{https://orcid.org/0009-0004-6753-4020}{0009-0004-6753-4020}
ライセンス：CC BY 4.0

\end{document}
